\scanpage{246}
. Nous allons expliciter maintenant la donnée de $Z$, en introduisant le sous groupe de $H$

\[
\deleted{\unsure{D}}\ T_H\overset{\mathrm{def}}{=}\{\,\deleted{\unsure{t}}u\in H\mid\exists\lambda\in Z,\ \deleted{\unsure{t}}us_0=\lambda s_0,\ ut_0=\lambda^{-1}t_0\}
\tag*{(\unsure{80|88})}
\]

C'est bien un sous-groupe, \deleted{puisque si $u$ et} si $u$ correspond à $\lambda$, $v$ correspond à $\mu$, alors $vu$ correspond à $\mu\lambda$. \deleted{D'ailleurs le $\lambda$ associé} à $u$ est unique, car \deleted{\unsure{on est ramené}} \inserted{$Z$ opère \uline{librement} sur $S$} \deleted{voir que} \deleted{mieux}

\[
\deleted{(\lambda s_0,\lambda^{-1}t_0)=(s_0,t_0)\Longrightarrow\lambda=1\quad(\lambda\in Z)}
\]

\deleted{ce qui résulte aussitôt de (72).} On a donc un hom de groupes canonique $T_H\longrightarrow Z$, dont le noyau est $\unsure{\overline{\mathfrak z}|\overline\zeta|\mathfrak z_H}$,

\diagram{\begin{tikzcd}[ampersand replacement=\&]
{(81)}\&1\arrow[r,no head]\&\mathfrak z_H\arrow[r,arrows={-|}]\&T_H\arrow[r,"{\unsure{\kappa|\pi|k}}"]\&Z\arrow[r]\&1
\end{tikzcd}}

\margin{Attention, cet hom \uline{n'est pas} l'hom $\delta\mid Z$}

qui dans le cas simplement transitif s'écrit

\[
T_H\xrightarrow{\sim}Z\quad .
\tag*{$(81')$}
\]

Notons

\[
T_H\text{ normalise }U_0,\ \inserted{\text{et }U_1,}\quad T_H\cap U_0=T_H\cap U_1=\mathfrak z_H
\tag*{(82)}
\]

\[
\begin{aligned}
&\deleted{t\in T_H\Longrightarrow\sigma(t)\ \unsure{=}\ t^{-1}\quad\text{i.e. }\sigma\text{ normalise }T_H}\\
&\deleted{\text{et induit sur }T_H/\unsure{\mathfrak z_H}\ \unsure{\simeq}\ \text{l'aut}}\\
&\deleted{\text{du passage à l'inverse}}
\end{aligned}
\]

En effet, si $\deleted{\unsure{u}}t\in T_H$, $u\in U_0$, on a

\[
tut^{-1}(s_0)=t\underbrace{u\,\lambda^{-1}(s_0)}_{\lambda^{-1}\underbrace{u(s_0)}_{s_0}}=t(\lambda^{-1}s_0)=\lambda^{-1}ts_0=\lambda^{-1}\lambda s_0=s_0
\]

d'où $tut^{-1}\in U_0$, \deleted{\unsure{de même on}} \inserted{par symétrie on} voit \deleted{\unsure{tout}} que $T_H$ normalise $U_1$. Si $t\in T_H\cap U_0$, on a

\[
t(s_0,t_0)=(\lambda s_0,\lambda^{-1}t_0)=(s_0,\lambda^{-1}t_0)\qquad\text{puisque }ts_0=s_0
\]

% archive sheet number: 245
